\chapter{Lorenz Curves, Majorization and Concentration}
\label{ch:lorenz}

The Lorenz identity $Z = \int_0^1 \Lambda(s) \, ds$ is the hinge of the book, and the Lorenz curve has a history and a theory of its own that bear on how the identity can be used. This chapter sets out that theory: the origin of the curve in the measurement of inequality, its relation to majorization and the Hardy Littlewood P\'olya theorem, the concentration function\index{concentration function} of a likelihood ratio, and the sense in which the quantile function is a complete invariant of a random variable up to rearrangement.

\section{Lorenz's curve}

Max Lorenz\index{Lorenz curve} proposed his curve in 1905 as a way of displaying the concentration of wealth. Order the population by income from poorest to richest, and plot against the fraction $u$ of the population the fraction of total income that this poorest fraction receives. The result is a convex curve from $(0,0)$ to $(1,1)$ lying below the diagonal; the diagonal is perfect equality and the curve hugging the horizontal axis is perfect inequality. The area between the curve and the diagonal, doubled, is Gini's coefficient of 1912.

In the notation of Chapter~\ref{ch:vertical}, let $Y \ge 0$ be a random variable with distribution function $F$ and quantile function\index{quantile function} $F^{-1}(u) = \inf\{ y : F(y) \ge u \}$. The Lorenz curve is
\[
\cL_0(u) = \frac{1}{\EE Y} \int_0^u F^{-1}(v) \, dv , \qquad u \in [0, 1] .
\]
Our $\Lambda$ is the reversed quantile, $\Lambda(s) = F^{-1}(1 - s)$, ordering from richest to poorest, which is natural when the large values are what one is after, and $Z = \EE Y = \int_0^1 \Lambda$ is Lorenz's normalisation. The curve used in the rest of this book, and in the Notation, is the richest first curve $\cL(u) = Z^{-1} \int_0^u \Lambda$, which is concave and related to Lorenz's poorest first curve $\cL_0$ by $\cL(u) = 1 - \cL_0(1 - u)$. The curve itself is scale free; the integral $Z$ is the scale. Nested sampling computes $Z$; the shape of $\Lambda$, whose normalised integral is the Lorenz curve, is what the walk reveals along the way.

\section{Majorization and Hardy, Littlewood and P\'olya}

Two vectors $a, b \in \mathbb{R}^n$ with the same sum are compared by majorization\index{majorization}: $a \prec b$ if the partial sums of the decreasing rearrangement of $a$ are dominated by those of $b$, $\sum_{i \le k} a_{[i]} \le \sum_{i \le k} b_{[i]}$ for all $k$. The vector $b$ is more spread out than $a$. Hardy, Littlewood and P\'olya proved in 1929 that $a \prec b$ if and only if $a = Mb$ for a doubly stochastic matrix $M$, and if and only if $\sum \phi(a_i) \le \sum \phi(b_i)$ for every convex $\phi$. In Lorenz terms, $a \prec b$ if and only if the Lorenz curve $\cL_0$ of $a$ lies above that of $b$, equivalently the curve $\cL$ of $a$ lies below that of $b$.

The relevance to counting is the following. The density of states $N(\cdot)$ is a vector indexed by energy, with sum $|\cX|$, and the Boltzmann\index{Boltzmann distribution} weights $e^{-s/\gamma}$ tilt it. As $\gamma$ decreases the Boltzmann vector $(\pi_\gamma(x))_{x \in \cX}$ becomes more concentrated on low energies in the majorization order (the mass of the $k$ lowest energy states is nondecreasing as $\gamma$ decreases), and the Lorenz curve of the likelihood ordinate rises. The annealing\index{simulated annealing} schedule is a path in the majorization order from the uniform distribution on $\cX$ to the uniform distribution on the ground states, and the ratios $Z(\gamma_k) / Z(\gamma_{k-1})$ are the costs of the steps along that path. The free energy $\log Z(\gamma)$ is a convex function of $1/\gamma$, and its second derivative, which is the variance of the energy, is the local cost of annealing.

\section{The rearrangement invariance}

A random variable and its quantile function determine each other: $Y \overset{D}{=} F^{-1}(U)$ for $U$ uniform. Everything about the distribution of $Y$ is in $F^{-1}$, and nothing about the underlying space $\cX$ or the map $L$ survives except through $F^{-1}$. Two likelihoods $L_1$ and $L_2$ on different spaces with the same quantile function have the same $Z$, the same Lorenz curve, and the same vertical walk. This is the sense in which the vertical approach is a dimension reduction: the walk cannot tell a problem on $\{0,1\}^{1000}$ from a problem on $[0,1]$ if their quantile functions agree.

What the walk can tell is how hard it is to produce conditioned draws, and that depends on $\cX$ and $L$ in full. The rearrangement invariance isolates the part of the problem that is independent of the geometry of $\cX$, namely the staircase, from the part that depends on it, namely the level set sampler\index{level set sampler}s. For complexity theory the first part is where the problem is stated and the second is where its hardness lives. The vertical view does not dissolve the hardness; it says that the hardness is entirely in the samplers, and that nothing about the staircase itself, which is a monotone function on an interval, is hard.

\section{The concentration function}

Let $P$ and $Q$ be probability measures on $\cX$ with $P \ll Q$ and likelihood ratio $\ell = dP/dQ$. The concentration function of $P$ with respect to $Q$, introduced by Cifarelli and Regazzini in 1987 as a generalisation of the Lorenz curve, is
\[
\phi_P(u) = \sup\{ P(A) : Q(A) \le u \} = \int_0^u m^{-1}(1 - z) \, dz ,
\]
for $Q$ nonatomic, or with $A$ replaced by randomised tests $0 \le \psi \le 1$ when $Q$ has atoms, where $m$ is the distribution function of $\ell$ under $Q$ and $m^{-1}$ its quantile; the second expression is the Neyman Pearson lemma, since the most probable set of $Q$ measure $u$ under $P$ is the upper level set of the likelihood ratio. The function $\phi_P$ is concave, increasing, with $\phi_P(0) = 0$ and $\phi_P(1) = 1$, and $\phi_P(u) = u$ for all $u$ if and only if $P = Q$. Its departure from the diagonal measures the disagreement of $P$ and $Q$, and the total variation distance is $\sup_u (\phi_P(u) - u)$.

The concentration function is the Lorenz curve of the likelihood ratio and so it is the object that vertical likelihood\index{vertical likelihood} computes when $L = \ell$. In Bayesian terms, with $Q$ the prior and $L$ the likelihood, $P$ is the posterior and $\phi_P$ records how much posterior mass can be packed into a prior set of given size; the evidence $Z = \int L \, dQ$ is the normalisation. In counting terms, with $Q$ uniform on $\cX$ and $L = \Ind_{\{S = 0\}}$, the concentration function, with randomised tests as above, is $\phi_P(u) = \min\{1, u/\rho\}$, a straight line of slope $1/\rho$ up to $u = \rho$ and then flat, and the rare event probability $\rho$ is the point where it reaches one. Estimating $\rho$ is locating the kink.

\section{Nested concentration functions}

The product estimator\index{product estimator} has a clean description in terms of concentration functions. Let $Q_t = Q(\cdot \mid S \le m_t)$ be the conditioned measures of the ladder, with levels $0 = m_0 < m_1 < \dots < m_T$ and $Z_{m_T} = 1$. Each $Q_{t-1}$ is absolutely continuous with respect to $Q_t$ with likelihood ratio $\Ind_{\{S \le m_{t-1}\}} / \theta_t$ where $\theta_t = Z_{m_{t-1}} / Z_{m_t}$, and the concentration function of $Q_{t-1}$ with respect to $Q_t$ is $\phi_t(u) = \min\{1, u / \theta_t\}$. The ladder is a chain of concentration functions, each a kink at $\theta_t$, and the composition of the kinks is the kink of the whole, at $\rho = \prod \theta_t$. The vertical walk estimates the kinks one at a time; the identity $\rho = \prod \theta_t$ is the statement that concentration functions of nested conditionings compose multiplicatively at the kink.

\section{Monotonicity, concavity and the direction of the walk}

Since $\Lambda$ is nonincreasing, the Lorenz curve $\cL(u) = Z^{-1} \int_0^u \Lambda$ is concave. Monotonicity is what makes the walk from $s = 1$ to $s = 0$ well behaved: the increments $\int_{s_{k}}^{s_{k-1}} \Lambda$ on a geometric grid $s_k = e^{-k/n}$ are each at most $\Lambda(s_k)(s_{k-1} - s_k)$ and at least $\Lambda(s_{k-1})(s_{k-1} - s_k)$, and the difference between the Riemann sums with left and right endpoints is at most a telescoping sum, and so is bounded by $1/n$ times the total variation of $\Lambda$, hence by $\Lambda(0)/n$. This is a deterministic bracket of order $1/n$; the $O(n^{-2})$ rate of Chapter~\ref{ch:rates} is a different statement, a mean squared error for random nodes, and does not follow from it. Concavity also explains why the walk goes from high to low energy and not the reverse: the mass of $Z$ is near $s = 0$, where $\Lambda$ is large, and a walk that started at $s = 0$ would need to begin with a draw from the rarest level set, which is the problem itself. Starting at $s = 1$, where draws are free, and descending, is the only direction in which each step is prepared by the one before.
